\clearpage

\section {Systematic effects in the prediction of the single $\mu$-like ring event energy spectrum}
\label{sec:systematics}

%....................................................................................................................
%....................................................................................................................

\subsection {Input systematic parameters}
\label{sec:systematics:inputs}

This analysis considers \NSyst{} flux, cross section and Super-K detector systematic parameters 
plus 7 oscillation parameters, which can be treated as either systematic parameters or
as parameters of interest, depending on the analysis that is performed.

\begin{enumerate}

\item 
{\bf Flux and cross-section parameters} (65 parameters).\\

In this analysis, all the flux and cross-section systematic parameters are included in the ND280 data (BANFF) fit.
However, since ND280 is not sensitive to the electron neutrino cross section and the NC1$\gamma$ parameters they remain unconstrained
The parameters fit are 
11 $\nu_{\mu}$, 5 $\bar{\nu}_{\mu}$, 7 $\nu_{e}$ and 2 $\bar{\nu}_{e}$ neutrino mode (FHC) flux parameters, 
5 $\nu_{\mu}$, 11 $\bar{\nu}_{\mu}$, 2 $\nu_{e}$ and 7 $\bar{\nu}_{e}$ antineutrino mode (RHC) flux parameters 
and 15 cross-section parameters. 
Best-fit values of these parameters and a covariance matrix giving their correlated uncertainties are obtained by the BANFF group from a fit 
of the $\bar{\nu}_\mu$ and $\nu_\mu$ 
%CC0$\pi$-like, CC1$\pi$-like and CC other-like 
ND280 samples as described in Ref~\cite{BANFF2015}. 
A list of the BANFF parameters, their best-fit values and their prefit and postfit errors,
taken from TN-230 \cite{BANFF2016},
is given in 
\cref{tab:banff_uncertainties_flux_fhc},~\cref{tab:banff_uncertainties_flux_rhc},
\cref{tab:banff_uncertainties_xsec}.

The best-fit values are used to reweight the templates as described in \cref{sec:banff_weights}, and the
postfit errors in the covariance matrix are used as prior uncertainties in this analysis.  The Super-K
systematic parameters are shown in \cref{tab:banff_uncertainties_skdetfsi}.  All the cross section parameters
set to 1 correspond to the NEUT nominal value, except ``CC other shape'' for which the NEUT nominal value is
0.  For all the normalisation parameters, 1 is the nominal pre-banff fit value.

The effects of these correlated uncertainties are evaluated using the Cholesky method, that allows throws of
the systematic parameters according to their multivariate distribution taking into account existing
correlations.
%The covariance matrix is ``Cholesky decomposed'' 
%into the product of a lower triangular matrix (all elements above the diagonal are zero) and its transpose (in which all elements below the 
%diagonal are zero). Random tweaks of each parameter are independently drawn from a Gaussian distribution with mean zero and standard deviation 1, 
%and these random numbers are put into a vector. This vector is then multiplied by the lower triangular matrix to give a vector of tweaked 
%uncertainties that have the correct correlations between them. The BANFF parameters are then varied by these tweaked uncertainties, with all 
%parameters being varied simultaneously. This process may be repeated as many times as necessary. 

Some of the cross-section parameters have complicated, often non-linear, effects on true and reconstructed
neutrino energy bin normalisations; the effects of uncertainties in these parameters are evaluated by
reweighting the MC templates. Weights are calculated for tweaks from -3$\sigma$ to +3$\sigma$ in steps of
0.5$\sigma$, and interpolations are made between those weights using cubic splines (for MaQE only, weights are
calculated from -7$\sigma$ to +3$\sigma$). Cubic splines are made for the following cross-section
uncertainties:
\begin{itemize}
\item CCQE axial-mass scaling factor (MaQE)
\item Fermi momentum for $^{16}$O
\item Binding energy for $^{16}$O
\item Nucleon to $\Delta$ transition axial form factor (C$^{A}_{5}$) 
\item Resonance-production axial-mass scaling factor (MaRes)
\item Scale of isospin 1/2 nonresonant background
\item CC other shape
\end{itemize}
The splines are made in bins of theta, true and reconstructed energy, and the binning is given in Ref.\
\cite{splineinfo}. The $\theta$ binning is the same as for the analysis templates, shown in
\cref{sec:mc_templates:template_binning}.

The effects of the flux and cross-section normalisation uncertainties are also evaluated by reweighting the MC
templates, but no splines are made for them since the weight is identical to the tweak for these uncertainties
(each flux parameter is a normalisation for a limited range of true energy).  The cross section parameters
which simply normalise the overall spectrum component are:
\begin{itemize}
\item 2p-2h, applied to all the 2p-2h events (both neutrino and antineutrino)
\item 2p-2h-$\bar{\nu}$, applied to all the antineutrino 2p-2h events
\item CC coherent for $^{16}$O normalisation
\item NC coherent normalisation
\item NC other normalisation
\item NC $1\gamma$ normalisation
\item CC $\nu_{e}$ normalisation
\item CC $\bar{\nu}_{e}$ normalisation
\end{itemize}

Uncertainties due to ``second class currents" are implemented as an additional uncertainty on CC
$\bar{\nu}_{e}$ and CC $\nu_{e}$ parameters with a 2\% anti-correlated constraint between them.  The BANFF
parameters are applied to the MC templates as specified in appendix~\ref{app:systematics} in
\cref{tab:apply_syst_to_template_mulike_fhc,%
  tab:apply_syst_to_template_elike_fhc,%
  tab:apply_syst_to_template_mulike_rhc,%
  tab:apply_syst_to_template_elike_rhc}
%      tab:apply_syst_to_template_nueCConepi_fhc }.
The normalisation parameter of the NC$1\gamma$ component was not used in the previous analyses,
since it was included in NC other.

In the BANFF fit a bias on the best-fit flux parameters is present due to the degeneracy between the neutrino
flux and ND280 detector parameters.  For this reason a covariance matrix that accounts for the bias of the
flux as well as all the other parameters is summed in quadrature with the total post-fit covariance. Furthermore
the best-fit values have been corrected in order to remove the bias. A full description of the procedure is
given in \cite{BANFF2015}.  As shown in \cite{OurRun1617b17cJointAnalysis3Flavour} the impact of the bias on the
systematic parameters fitted by the BANFF analysis is small compared to the other systematic uncertainties.


%....................................................................................................................

\subsection {Comparison of predicted number of events in neutrino mode $e$-like sample between current and 2015 analysis}
\label{sec:systematics:comparison_events_elike_2015}

The predicted number of events in antineutrino mode $e$-like sample is now smaller compared to the
$\bar{\nu}_e$ appearance analysis presented at the 2015 summer conferences
\cite{OurRun56AntiNueBarAppearanceAnalysis3Flavour}, 3.218 instead of 3.729.  The differences between the
current and previous oscillation analyses are mainly due to:

\begin{enumerate}

\item new value of $\sin^2 \theta_{13}$ corresponding to the PDG 2015 release, as shown in
  \cref{tab:osc_params_joint_asimov}, which is 0.0217 while in the previous analysis it was 0.0250. If the old
  PDG value is used the expected number of antineutrino mode $e$-like events becomes 3.542.

\item Subsequently, if the cross section parameters are set to the old post-BANFF fit values, the expected
  number of events becomes 3.63. Indeed all the post-BANFF 2015 fit cross section parameters are higher
  compared to the 2016 BANFF fit result. The effect of the flux parameters is smaller, indeed if they are set
  to BANFF 2015 result, the expected number of events becomes 3.66.
 
\item Finally, after all the changes listed above, if the rescaling of CC coherent and NC$1\gamma$ components
  (see \cref{sec:mc_templates:input_mc}) is not applied, agreement between the 2015 and 2016 analysis is
  obtained. The first one increases the number of events up to 3.772 while the latter one reduces down to
  3.711, which is approximately the number of events expected in the 2015 analysis.

%\item the total normalisation for 2p-2h events for antineutrinos is 0.894,   �  
%	where the BANFF best-fit of the 2p-2h parameter applied only to $\bar{\nu}$ is $0.578$ and 
%	the best-fit of the 2p-2h parameter applied to all the flavors is 1.546.
%	 This reduces the expected total number 
%	of $\bar{\nu}_e$ events, while in the past analysis a single systematic parameter was applied to all the 2p-2h events 
%	and its BANFF best-fit was about 1.031.

\end{enumerate}


{
\renewcommand{\arraystretch}{1.2}
\begin{table}
\centering
\begin{tabular}{| c | c | c | c |}
\hline
           Parameter & Description & Best fit & 1$\sigma$ pre/postfit \\
           & & & fractional error \\
\hline
            0 &  $\nu_{\mu}$ flux normalisation, E = 0.0 - 0.7 GeV               & 1.142   & 0.094 / 0.071 \\ 
            1 &  $\nu_{\mu}$ flux normalisation, E = 0.7 - 1.0 GeV               & 1.091   & 0.079 / 0.053 \\ 
            2 &  $\nu_{\mu}$ flux normalisation, E = 1.0 - 1.5 GeV               & 1.079   & 0.077 / 0.052 \\ 
            3 &  $\nu_{\mu}$ flux normalisation, E = 1.5 - 2.5 GeV               & 1.074   & 0.081 / 0.060 \\ 
            4 &  $\nu_{\mu}$ flux normalisation, E = 2.5 - 30.0 GeV              & 1.050   & 0.080 / 0.053 \\ 
            5 &  $\bar{\nu}_{\mu}$ flux normalisation, E = 0.0 - 0.4 GeV         & 1.204   & 0.104 / 0.069 \\ 
            6 &  $\bar{\nu}_{\mu}$ flux normalisation, E = 0.4 - 0.5 GeV         & 1.210   & 0.102 / 0.058 \\ 
            7 &  $\bar{\nu}_{\mu}$ flux normalisation, E = 0.5 - 0.6 GeV         & 1.180   & 0.096 / 0.051 \\ 
            8 &  $\bar{\nu}_{\mu}$ flux normalisation, E = 0.6 - 0.7 GeV         & 1.128   & 0.085 / 0.041 \\ 
            9 &  $\bar{\nu}_{\mu}$ flux normalisation, E = 0.7 - 1.0 GeV         & 1.123   & 0.125 / 0.077 \\ 
           10 &  $\bar{\nu}_{\mu}$ flux normalisation, E = 1.0 - 1.5 GeV         & 1.118   & 0.105 / 0.068 \\ 
           11 &  $\bar{\nu}_{\mu}$ flux normalisation, E = 1.5 - 2.5 GeV         & 1.121   & 0.080 / 0.053 \\ 
           12 &  $\bar{\nu}_{\mu}$ flux normalisation, E = 2.5 - 3.5 GeV         & 1.121   & 0.074 / 0.049 \\ 
           13 &  $\bar{\nu}_{\mu}$ flux normalisation, E = 3.5 - 5.0 GeV         & 1.122   & 0.094 / 0.067 \\ 
           14 &  $\bar{\nu}_{\mu}$ flux normalisation, E = 5.0 - 7.0 GeV         & 1.112   & 0.093 / 0.064 \\ 
           15 &  $\bar{\nu}_{\mu}$ flux normalisation, E = 7.0 - 30.0 GeV        & 1.093   & 0.130 / 0.108 \\ 
           16 &  $\nu_{e}$ flux normalisation, E = 0.0 - 2.5 GeV                 & 1.109   & 0.069 / 0.047 \\ 
           17 &  $\nu_{e}$ flux normalisation, E = 2.5 - 30.0 GeV                & 1.090   & 0.085 / 0.067 \\ 
           18 &  $\bar{\nu}_{e}$ flux normalisation, E = 0.0 - 0.5 GeV           & 1.193   & 0.095 / 0.054 \\ 
           19 &  $\bar{\nu}_{e}$ flux normalisation, E = 0.5 - 0.7 GeV           & 1.184   & 0.091 / 0.047 \\ 
           20 &  $\bar{\nu}_{e}$ flux normalisation, E = 0.7 - 0.8 GeV           & 1.167   & 0.091 / 0.049 \\ 
           21 &  $\bar{\nu}_{e}$ flux normalisation, E = 0.8 - 1.5 GeV           & 1.141   & 0.084 / 0.044 \\ 
           22 &  $\bar{\nu}_{e}$ flux normalisation, E = 1.5 - 2.5 GeV           & 1.120   & 0.080 / 0.054 \\ 
           23 &  $\bar{\nu}_{e}$ flux normalisation, E = 2.5 - 4.0 GeV           & 1.123   & 0.089 / 0.068 \\ 
           24 &  $\bar{\nu}_{e}$ flux normalisation, E = 4.0 - 30.0 GeV          & 1.180   & 0.156 / 0.144 \\ 
           25 &  Meson exchange current normalisation for $^{16}$O               & 1.031   & 1.041 / 1.015 \\ 
           26 &  C$^{A}_{5}$ nucleon to $\Delta$ transition axial form factor    & 0.854   & 0.119 / 0.073 \\ 
           27 &  Scale of isospin 1/2 nonresonant background                     & 1.143   & 0.154 / 0.146 \\ 
           28 &  CCQE axial-mass scaling factor                                  & 0.948   & 0.058 / 0.028 \\ 
           29 &  Resonance-production axial-mass scaling factor                  & 0.762   & 0.158 / 0.055 \\ 
           30 &  Fermi momentum for $^{16}$O                                     & 0.997   & 0.055 / 0.054 \\ 
           31 &  CC other shape                                                  & 0.023   & 0.400 / 0.193 \\ 
           32 &  Binding energy for $^{16}$O                                     & 1.002   & 0.333 / 0.326 \\ 
           33 &  CC coherent for $^{16}$O normalisation                          & 1.076   & 1.000 / 0.972 \\ 
           34 &  NC coherent normalisation                                       & 0.980   & 0.300 / 0.299 \\ 
           35 &  NC other normalisation                                          & 1.413   & 0.300 / 0.186 \\ 
           36 &  CC $\nu_{e}$ normalisation                                      & 1.001   & 0.020 / 0.020 \\ 
\hline
\end{tabular}
\caption{
 Summary of flux and correlated cross-section (BANFF) systematics included in the {\em VALOR} Run 5+6 
 $\bar{\nu}_{\mu}$-disappearance analysis.
 A description of how the BANFF fit constrained these systematics can be found in Ref.~\cite{BANFF2015}.
}
\label{tab:banff_uncertainties}
\end{table}
}

\clearpage


\item
{\bf Super-K efficiencies \& effects of intranuclear and secondary re-interactions} (44 parameters).\\

The Super-K detector uncertainties include the efficiencies of the fiducial volume and reduction chain, and
the OD, ring-counting, PID, momentum and decay-electron cuts. These uncertainties have been evaluated by
comparing 1632.2 days of SK-IV atmospheric data with atmospheric neutrino Monte Carlo. Ring counting errors
are evaluated using $\nu_\mu$ CCQE, $\nu_\mu$ CCnQE, NC, and $\nu_e$ CC-enriched control samples from
atmospheric neutrinos.  The NC PID error is evaluated using a NC-enriched control sample, obtained by not
using the PID information during the selection. All the Monte Carlo control samples are fitted simultaneously
to the data with a $\chi^2$ function using the selection efficiencies as fit parameters
~\cite{SuperkErrors_numu}~\cite{SuperkErrors_nue}.

Intranuclear final-state interactions (FSI) have significant effects in the T2K energy range and result in
event topologies, for scattering from nuclear targets, which are drastically different from those created in
neutrino scattering from free nucleons.  Uncertainties due to FSI were estimated by simultaneously varying the
NEUT parameters that scale the interaction probabilities for quasi-elastic scattering, absorption, charge
exchange and pion production using the method described in Sec.\ 2 in \cite{NIWGErr}. The NIWG group produced
16 parameter sets that span the 1$\sigma$ surface in the parameter space; this was done by varying the NEUT
parameters when running the pion scattering simulations, and comparing the results with external pion-carbon
scattering data. These 16 parameter sets are shown in Tab.\ 2 in \cite{NIWGErr}.  Uncertainties due to
secondary pion interactions (SI) were also evaluated. The NIWG group adapted the NEUT nuclear cascade model to
simulate SI in the Super-K detector simulation (SKDETSIM) and to allow variations in the interaction
probabilities. Details of this are given in \cite{SI}.  Since the same model is used for FSI and SI, it is
possible to evaluate the uncertainties in both FSI and SI simultaneously. This was done by reweighting a
single $e$-like or $\mu$-like ring selected SK MC with the 16 FSI parameter sets using T2KReweight. The
reweighting method for FSI is described in Sec.\ 8 in \cite{NEUTnuclear}. Variations in the SI interaction
probabilities were included for each FSI parameter set. Subsequently a covariance matrix was calculated using
equation 1 in \cite{NIWGErr}; the quantity p in that equation is the average event weight for a FSI parameter
set, and the sum is over the parameter sets.


\begin{table}
\centering
\begin{tabular}{| c | c | c | r |}
\hline
 Parameter & Sample        & Category                              & Energy range (GeV)                  \\
\hline
 0 ~$|$ 18 & electron-like & oscillated $\nu_e$/$\bar{\nu}_e$ CC   & $E_{reco}<$ 0.35                      \\ 
 1 ~$|$ 19 & electron-like & oscillated $\nu_e$/$\bar{\nu}_e$ CC   & 0.35 $<E_{reco}<$ 0.80                \\ 
 2 ~$|$ 20 & electron-like & oscillated $\nu_e$/$\bar{\nu}_e$ CC   & 0.80 $<E_{reco}<$ 1.25                \\ 
 3 ~$|$ 21 & electron-like & $\nu_\mu$/$\bar{\nu}_{\mu}$ CC         & $E_{reco}<$ 0.35                      \\ 
 4 ~$|$ 22 & electron-like & $\nu_\mu$/$\bar{\nu}_{\mu}$ CC         & 0.35 $<E_{reco}<$ 0.80                \\ 
 5 ~$|$ 23 & electron-like & $\nu_\mu$/$\bar{\nu}_{\mu}$ CC         & 0.80 $<E_{reco}<$ 1.25                \\ 
 6 ~$|$ 24 & electron-like & $\nu_e$/$\bar{\nu}_{e}$ CC            & $E_{reco}<$ 0.35                       \\ 
 7 ~$|$ 25 & electron-like & $\nu_e$/$\bar{\nu}_{e}$ CC            & 0.35 $<E_{reco}<$ 0.80                 \\ 
 8 ~$|$ 26 & electron-like & $\nu_e$/$\bar{\nu}_{e}$ CC            & 0.80 $<E_{reco}<$ 1.25                 \\ 
 9 ~$|$ 27 & electron-like & all NC                                & $E_{reco}<$ 0.35                       \\ 
 10 $|$ 28 & electron-like & all NC                                & 0.35 $<E_{reco}<$ 0.80                 \\ 
 11 $|$ 29 & electron-like & all NC                                & 0.80 $<E_{reco}<$ 1.25                 \\ 
\hline
 12 $|$ 30 & muon-like     & $\nu_\mu$/$\bar{\nu}_{\mu}$ CCQE       & $E_{reco}<$ 0.40                       \\ 
 13 $|$ 31 & muon-like     & $\nu_\mu$/$\bar{\nu}_{\mu}$ CCQE       & 0.40 $<E_{reco}<$ 1.10                 \\ 
 14 $|$ 32 & muon-like     & $\nu_\mu$/$\bar{\nu}_{\mu}$ CCQE       & 1.10 $<E_{reco}<$ 30.0                 \\ 
 15 $|$ 33 & muon-like     & all $\nu_\mu$/$\bar{\nu}_{\mu}$ CCnQE  & $E_{reco}<$ 30.0                       \\ 
 16 $|$ 34 & muon-like     & $\nu_e$/$\bar{\nu}_{e}$ CC             & $E_{reco}<$ 30.0                      \\ 
 17 $|$ 35 & muon-like     & all NC                                & $E_{reco}<$ 30.0                       \\ 
\hline
 36~$|$ --  & \nueCConepi{}-like & oscillated $\nu_e$/$\bar{\nu}_e$ CC   & 0.35 $<E_{reco}<$ 0.80            \\ 
 37~$|$ --  & \nueCConepi{}-like & oscillated $\nu_e$/$\bar{\nu}_e$ CC   & 0.80 $<E_{reco}<$ 1.25            \\ 
 38~$|$ --  & \nueCConepi{}-like & $\nu_\mu$/$\bar{\nu}_{\mu}$ CC         & 0.35 $<E_{reco}<$ 0.80            \\ 
 39~$|$ --  & \nueCConepi{}-like & $\nu_\mu$/$\bar{\nu}_{\mu}$ CC         & 0.80 $<E_{reco}<$ 1.25            \\ 
 40~$|$ --  & \nueCConepi{}-like & $\nu_e$/$\bar{\nu}_{e}$ CC            & 0.35 $<E_{reco}<$ 0.80             \\ 
 41~$|$ --  & \nueCConepi{}-like & $\nu_e$/$\bar{\nu}_{e}$ CC             & 0.80 $<E_{reco}<$ 1.25             \\ 
 42~$|$ --  & \nueCConepi{}-like & all NC                                & 0.35 $<E_{reco}<$ 0.80             \\ 
 43~$|$ --  & \nueCConepi{}-like & all NC                                & 0.80 $<E_{reco}<$ 1.25             \\         
\hline
\end{tabular}
\caption{
Binning of Super-K selection efficiency 
%and FSI+SI 
systematic parameters in covariance matrices for electron-, muon-, and \nueCConepi{}-like  samples
for each category and reconstructed neutrino energy range. The same parametrization is used for both neutrino and antineutrino mode samples.
The parameter number is shown for both neutrino (0-17) and antineutrino (18-35) mode for the 0\( \pi \) samples, and neutrino (36-43) mode for the \nueCConepi{} sample.
}
\label{tab:skdet_params}
\end{table}



We consider the Super-K detector efficiency and FSI+SI systematics together, and add their covariance matrices
linearly. Each of these covariance matrices has \textcolor{red}{\textbf{SET parameter numbers}} correlated
parameters, corresponding to different modes and reconstructed energy bins, as shown in
\cref{tab:skdet_params}. The effects of these uncertainties are also evaluated using the Cholesky method. \\

\item
{\bf Super-K energy scale} (1 parameter).\\

The systematic parameter $f_{E; r}^{SK}$ is included to estimate the effects of uncertainty in the Super-K
reconstructed energy scale.  This uncertainty is estimated to be 2.4\% \cite{SuperkEScale2013}; its effects
are calculated by scaling the bin edges of the MC templates and, assuming uniform distribution of events
within the bins, calculating the number of events gained from / lost to neighbouring bins.  The effect of the
energy scale uncertainty is propagated to all MC templates as shown in
\cref{tab:apply_syst_to_template_mulike_fhc,%
      tab:apply_syst_to_template_elike_fhc,%
      tab:apply_syst_to_template_mulike_rhc,%
      tab:apply_syst_to_template_elike_rhc}
%      tab:apply_syst_to_template_nuecc1pi_fhc}
and is common to all the Super-K samples. \\


\item 
{\bf Oscillation parameters and mass hierarchy} (7 parameters).\\

The 3-flavour oscillation probability is used.  All the nuisance oscillation parameters, i.e. those parameters
which are not shown in a given contour, are marginalised (see \cref{sec:oscfit}) except the mass hierarchy
which is fixed to either normal or inverted hierarchy.  The treatment of the oscillation parameters in the
analysis is shown in \cref{tab:osc_params_joint_priors} while the nominal values are shown in
\cref{tab:osc_params_joint_asimov}.  Two different sets of oscillation parameters were chosen. The Asimov data
set A is the best-fit of the Run 1-4 data Mach3 analysis \cite{Mach3JointFitRun14}. The Asimov data set B is a
variation of Asimov data set A, which is useful for testing the sensitivity to the opposite
$\sin^2 \theta_{23}$ octant, which reduces the number of expected events in the $e$-like samples, as well as
the sensitivity to $\delta_{CP}$.

%Differently from the $\bar{\nu}_{\mu} \rightarrow \bar{\nu}_e$ appearance analysis,
%where the prior distributions corresponded to the T2K Run 1-4 best-fit probability distributions
%of the Mach3 Bayesian analysis, now priors from the global fit of \cite{pdg_2014} are used.
For $\sin^2 \theta_{13}$ two different priors can be used, depending on whether the reactor constraint is
applied, denoted as ``reactors", when applied and ``T2K only" when not applied.  In the first case a gaussian
prior given by the combined measurement of reactor experiments is used, while in the latter case a flat prior
which covers a large range of values is adopted.  With respect to the older global fit \cite{pdg_2013}, the
value of $\sin^2 \theta_{13}$ has been reduced by about 10\% and the associated error is also reduced.  Since
an excess of events was seen in the T2K Run 1-4 neutrino mode $e$-like sample and not much data are collected
in neutrino mode during the Run 5+6 (only $0.34 \times 10^{20}$ POT), this change is expected to have an
impact on the search for CP violation, making the confidence intervals tighter. \textcolor{red}{\textbf{UPDATE with run 7a and 7c?}}


\begin{table}
\centering
 \begin{tabular}{| c | c | c |}
% \hline
% Parameter(s)                                               			&     Treatment              &           Nominal value (NH)           & Nominal value (IH)      \\
%\hline
%$\sin^{2}\theta_{23}$                                       			&     marginalized         &                \MixingTwoThreeNH       &    \MixingTwoThreeIH  \\
%$\sin^{2}\theta_{13}$  							&     marginalized         &                \MixingOneThreeNH	&      \MixingOneThreeIH	 \\
%$\sin^{2}\theta_{12}$  							&     fixed               	    &             \MixingOneTwoIH             &    \MixingOneTwoIH	 \\
%$|\Delta m^{2}_{32}|$ (NH) / $|\Delta m^{2}_{31}|$ (IH)      &     marginalized          &               \DmsqNH			&      \DmsqIH                        \\
%$\Delta m^{2}_{21}$      							&     fixed                       &             \dmsqIH                          & 	\dmsqIH		\\
%$\delta_{CP}$                                               			&     marginalized           &              \DeltaCPNH			&     \DeltaCPIH              \\
%Mass Hierarchy									&     marginalized	    &           Normal Hierarchy
 \hline
 Parameter(s)                                               			&     Treatment              &           Nominal value 		    \\
\hline
$\sin^{2}\theta_{23}$                                       			&     marginalized         &                \MixingTwoThreeNH        \\
$\sin^{2}\theta_{13}$  							&     marginalized         &                \MixingOneThreeNH	 \\
$\sin^{2}\theta_{12}$  							&     fixed               	    &             \MixingOneTwoIH              \\
$|\Delta m^{2}_{32}|$ (NH) / $|\Delta m^{2}_{31}|$ (IH)      &     marginalized          &               \DmsqNH		          \\
$\Delta m^{2}_{21}$      							&     fixed                       &             \dmsqIH                         	\\
$\delta_{CP}$                                               			&     marginalized           &              \DeltaCPNH		      \\
Mass Hierarchy									&     marginalized	    &            	NH         \\  
\hline
\end{tabular}
\caption{
Treatment of the oscillation parameters in the p-value computation as well as used values for fixed parameters.
The values of the fixed oscillation parameters correspond to the most probable parameters of the bayesian analysis
described in \cite{T2KLongPaper,Mach3JointFitRun14}.
The mass hierarchy is marginalized as well.
}
\label{tab:osc_params_fixed_values_reactor}
\end{table}



\end{enumerate}




%%%% SYSTEMATIC ERROR ENVELOPE PLOTS

%
%The effects of all the combined systematic uncertainties on the reconstructed energy spectrum of single ring events
%is shown in \cref{fig:spectra_systematic_Asimov_A} 
%for the typical oscillation scenario with parameters of Asimov datasets A.
%In appendix~\cref{app:systematics} the same plots are shown for different oscillation scenarios.
%These plots show the total error envelope for the reconstructed energy spectra; 
%this is calculated
%as the $\pm$1$\sigma$ spread with respect to the median of bin contents from \NumToys toy MC experiments generated with 
%randomized systematic parameters.
%All systematic parameters are considered and their correlations are taken into account.
%The spectra were tuned with the systematic parameters measured with the ND280 data 
%shown in \cref{tab:banff_uncertainties_flux_fhc},~\cref{tab:banff_uncertainties_flux_rhc}, 
%\cref{tab:banff_uncertainties_xsec},
%the nominal Super-K systematic parameters shown in \cref{tab:banff_uncertainties_skdetfsi}
%and the oscillation parameter priors shown in \cref{tab:osc_params_joint_priors}.
%In \cref{fig:Osc_envelope_Asimov_A},~\cref{fig:SKDet_FSI_envelope_Asimov_A},~\cref{fig:flux_envelope_Asimov_A},
%\cref{fig:xsec_envelope_Asimov_A},~\cref{fig:flux_xsec_envelope_Asimov_A} 
%the error envelopes and fractional errors of reconstructed energy spectra for all the selected samples,
%given respectively by
%oscillations, Super-K detector + FSI/SI, flux only, cross section only, flux+cross section systematics
%are shown. 
%
%The effects of systematic uncertainties on the predicted number of events are summarised in 
%\cref{tab:nsk_systematic_table_summary_joint_asimov_A_fhc_mulike},
%\cref{tab:nsk_systematic_table_summary_joint_asimov_A_fhc_elike},
%\cref{tab:nsk_systematic_table_summary_joint_asimov_A_rhc_mulike} and
%\cref{tab:nsk_systematic_table_summary_joint_asimov_A_rhc_elike} 
%\cref{tab:nsk_systematic_table_summary_joint_asimov_A_ccpi} 
%with the typical parameter values given in \cref{tab:osc_params_joint_asimov}.
%For the measurement of $\delta_{CP}$ it is important to understand the effect of the systematic uncertainties
%on the ratio of the number of events between neutrino and antineutrino mode samples,
%which is the consequence of correlations between systematic parameters.
%If correlations do not reproduce the physics correctly, a fake asymmetry between neutrino and antineutrino oscillations can be seen.
%On the other hand if correlations exist but are not taken into account, the measurement of $\delta_{CP}$ could be biased.
%Depending on the value of $\delta_{CP}$, correlations or anti-correlations between parameters are more conservative.
%If true $\delta_{CP} = \pm \pi/2$ the most conservative case is given by anti-correlation in the $e$-like samples, 
%since the oscillation asymmetry could be hidden by the systematic uncertainties.
%If true $\delta_{CP} \sim 0$ uncorrelated parameters reduce the sensitivity to $\delta_{CP}$.
%The effect of the systematic parameters on the ratio between neutrino and antineutrino mode events are shown in 
%\cref{tab:ratio_systematic_table_summary_asimov_A} and~\cref{tab:ratio_systematic_table_summary_asimov_B} in 
%appendix~\cref{app:systematics}.
%The most critical systematic parameters are $\sigma_{\nu_e}$ and $\sigma_{\bar{\nu}_e}$
%which are applied respectively to $\nu_e$ and $\bar{\nu}_e$ events.
%Their fractional error is 2.83\%
%and are partially anti-correlated with a non-diagonal elements in the covariance matrix of -0.0004 
%\cite{NIWG_reccomandation}.
%
%For the systematic effect study, in order to reduce the CPU time, the toys MC were produced using 
%%2-dimensional templates $E_{reco}$ vs $\theta$ for $e$-like samples
%%and the projection onto the $E_{reco}$ axis was performed
%1-dimensional templates $E_{reco}$ $e$-like templates,
%instead of the 2-dimensional $E_{reco}$ vs $\theta$,
%used in the analysis.
%However, as shown in appendix~\cref{app:comp_1d_ereco_2d_ereco_vs_theta},
%it has been checked that using 1- or 2- dimensional 
%templates does not make any difference in the 
%drawing of the error envelopes and the effect of the systematic
%uncertainties on the total number of events.
%


\clearpage


%..................


%%% Local Variables:
%%% mode: latex
%%% TeX-master: "jointfit_5sample2016"
%%% End:
